% Edited conversion of source pp. 32--33.
\ReportHeading
  {Особливості динамічної поведінки багатошарових систем, зумовлені їх шаруватою структурою та початковими напруженнями}
  {Ю.~П. Глухов}
  {Інститут механіки ім. С.~П. Тимошенка НАН України}
  {}

Шаруваті структури широко застосовують у транспорті, будівництві та авіації, оскільки вони дають змогу раціонально поєднувати матеріали й досягати кращих механічних характеристик, ніж в однорідних аналогах. Реальні конструкції часто працюють за наявності початкових, або залишкових, напружень і зазнають дії рухомих навантажень, що формує складну хвильову картину.

У межах лінеаризованої теорії пружності тіл із початковими напруженнями сформульовано просторові та плоскі задачі про реакцію багатошарової основи на поверхневе навантаження, яке рухається з постійною швидкістю. Розглянуто систему з плоско-паралельними межами на пружному півпросторі або жорсткій основі. Матеріали шарів описано нелінійно-пружними моделями з довільним пружним потенціалом, а додаткові напруження вважаються малими порівняно з початковими.

Розв’язки одержано за допомогою регуляційно-спектрального методу зі структурною регуляризацією. Підхід виявляє ефекти, що виникають унаслідок взаємодії шарів і не зводяться до суперпозиції властивостей окремих матеріалів.

\begin{enumerate}[label=\arabic*.,leftmargin=2em,itemsep=0.25em]
\item \textbf{Ефекти математичної структури моделі.}
  \begin{enumerate}[label*=\arabic*.,leftmargin=2.2em,itemsep=0.12em]
  \item \textit{Спектральна чутливість.} Залежно від початкових напружень і швидкості руху навантаження характеристичні рівняння можуть мати дійсні кратні корені навіть у дозвуковому режимі.
  \item \textit{Асиметрія хвильового поля.} У транс- і надзвукових режимах зміна типу рівнянь руху спричиняє конуси Маха та різне затухання прямих і зворотних хвиль; у шаруватій системі додається інтерференція на межах поділу.
  \item \textit{Локальність впливу початкових напружень.} Збурення від напруженого шару затухають із відстанню за степеневим або експоненціальним законом, тому домінують у ближній зоні.
  \end{enumerate}
\item \textbf{Ефекти фізичної природи.}
  \begin{enumerate}[label*=\arabic*.,leftmargin=2.2em,itemsep=0.12em]
  \item \textit{Системний вплив.} Попередньо напружений шар змінює ефективну жорсткість і умови поширення хвиль, перерозподіляючи напруження у сусідніх шарах.
  \item \textit{Асиметрія стиску та розтягу.} Попередній стиск підвищує ефективну жорсткість у межах лінеаризованої теорії, а розтяг її знижує, змінюючи амплітуди й затухання збурень.
  \item \textit{Зони слабкої чутливості.} За певного співвідношення довжин хвиль і глибини залягання напружених шарів початковими деформаціями можна нехтувати без істотної втрати точності.
  \item \textit{Жорсткість поверхневого шару та контакт.} Жорсткіший верхній шар екранує нижні, а жорстке з’єднання розсіює енергію рівномірніше, ніж ковзний контакт.
  \end{enumerate}
\end{enumerate}

Отримані результати показують, що динаміка багатошарової основи визначається одночасно математичною структурою крайової задачі та фізичними механізмами хвильової взаємодії, інтерференції й перетворення енергії на межах шарів.
